% ------------------------------------------------- %
% ---   NAGŁÓWEK - PROSZĘ NIE ZMIENIAĆ USTAWIEŃ --- %
% ------------------------------------------------- %
\documentclass[a4paper,10pt]{article}
\usepackage[utf8]{inputenc}
\usepackage[polski]{babel}
\usepackage[T1]{fontenc}
% Zamiast: \usepackage[margin=1.8cm]{geometry}
\usepackage[margin=1cm, top=1.4cm, headsep=0.3cm]{geometry}
\usepackage{tabularx}
\usepackage{booktabs}
\usepackage{fontawesome5}
\usepackage{amsmath,amssymb}
\usepackage[table]{xcolor}
\usepackage{graphicx}
\usepackage{fancyhdr}

\definecolor{tableheader}{RGB}{230, 238, 248}
\definecolor{scorebg}{RGB}{205, 222, 242}
\definecolor{primaryblue}{RGB}{30, 75, 140}

\pagestyle{fancy}
\fancyhf{}
\lhead{\small \textbf{Projekt 1} -- Wstęp do Uczenia Maszynowego 2026Z}
\rhead{\small Raport KM1}
\renewcommand{\headrulewidth}{0.4pt}

% ------------------------------------------------------ %
% --- KONIEC NAGŁÓWKA - PROSZĘ NIE ZMIENIAĆ USTAWIEŃ --- %
% ------------------------------------------------------ %

\begin{document}


\vspace{-1em}
\noindent
\renewcommand{\arraystretch}{1.3}
\begin{tabularx}{\textwidth}{|>{\hsize=1.2\hsize}X|>{\hsize=0.8\hsize}X|>{\centering\arraybackslash}m{2.5cm}|}
\hline
\rowcolor{tableheader}

% ----------------------------------------------------------- %
% ---   PROSZĘ UZUPEŁNIĆ TABELĘ W LINIJKACH Z KOMENTARZEM --- %
% ----------------------------------------------------------- %

\textbf{Imię i Nazwisko:} Paweł Buczyński % WPISZ IMIE I NAZWISKO
& \textbf{Kierunek:} IAD % ZOSTAW WŁAŚCIWE
& \cellcolor{scorebg}\textbf{PUNKTACJA} \\
\rowcolor{tableheader}
\textbf{Prowadzący:}  AK % ZOSTAW WŁAŚCIWE
& \textbf{Numer albumu:} 338979 % WPISZ NUMER ALBUMU 
& \cellcolor{scorebg}\huge{\textbf{ / 8}} \\
\hline
\end{tabularx}
\renewcommand{\arraystretch}{1.0}

\vspace{-1em}


\section*{\color{primaryblue}Raport z Kamienia Milowego 1: Eksploracja i przygotowanie danych}

% ----------------------- %
% ---  TREŚĆ RAPORTU  --- %
% ----------------------- %


\section{Eksploracyjna analiza danych}

Zbiór treningowy składa się z 6400 obserwacji i 23 zmiennych opisujących klientów kart kredytowych. Zmienne obejmują informacje demograficzne, historię spłat zadłużenia oraz wysokość rachunków i dokonanych wpłat.

\subsection{Analiza rozkładów}

Zmienna kategoryczna MARRIAGE przyjmuje w większości dwie wartości, ale jest też niewielka liczba obserwacji o innej, trzeciej wartości.
Być może jest to reprezentacja braku danych.
Pdobnie w zmiennej EDUCATION istnieje kategoria z bardzo małą liczbą obserwacji.

Zmienna AGE nie odzwierciedla typowej struktury demograficznej rozwiniętego społeczeństwa - widać tu znacznie mniejszy udział osób starszych.

Zmienne PAY\_x, określające terminowość spłaty, są skupione w zakresie 3-4 dni wokół zera i prawie symetryczne.
Jest nieco więcej dodatnich outlierów, ale wszystkie spłaty mieszczą się w zakresie około tygodnia od terminu.

Zmienne BILL\_AMTx, określające kwotę miesięcznego wyciągu, są wyraźnie prawostronnie skośne, co odzwierciedla typowy rozkład zamożności ludzi (podobnie jak LIMIT\_BAL). Występuje również kilka outlierów z dużą ujemną wartością.

Zmienne PAY\_AMTx, określające kwoty wpłacane na rzecz spłaty długu, są jeszcze bardziej skośne.
Większość wpłat mieści się w zakresie rzędu 5-10 tysięcy, ale są outliery rzędu setek tysięcy.
Być może są to spłaty całego zadłużenia na raz zamiast w ratach, albo po prostu klienci operujący wielkimi kwotami na codzień.

Występują braki danych w zmiennych PAY\_5, BILL\_AMT6 oraz PAY\_AMT4, stanowią one około 10\% obserwacji. Na tym etapie projektu je ignoruję.

\subsection{Analiza zależności}

Zmienne PAY\_x, określające terminowość spłaty, są ze sobą wyraźnie skorelowane; do tego korelacja między bliższymi miesiącami jest silniejsza -  ma to sens.

Zmienne BILL\_AMTx, określające kwotę miesięcznego wyciągu, są ze sobą silnie skorelowane.

Natomiast zmienne PAY\_AMTx, określające kwoty wpłacane na rzecz spłaty długu, są ze sobą dużo słabiej skorelowane.
Sugeruje to dużą nieregularność w spłacaniu długów, lub przeplatanie spłat z kolejnymi zadłużeniami/wydatkami.
Co ciekawe, widać wyraźnie, że wpłacana kwota jest najmocniej skorelowana z kwotą wyciągu w poprzednim miesiącu - co też ma sens.

Zmienna docelowa (y\_train) jest widocznie skorelowana z terminowością poprzednich wpłat, najsilniej z ostatnimi, a także nieco mniej z kwotami wpłat i limitem kredytowym.

\section{Metryka ewaluacji i podział danych}

Jako główną metrykę ewaluacji wybrałem ROC AUC. Pozwala ona ocenić zdolność modelu do rozróżniania klientów terminowo spłacających zadłużenie od klientów zagrożonych opóźnieniem, niezależnie od przyjętego progu klasyfikacji.
Zbiór danych nie jest mocno niezbalansowany, więc nie wymusza to użycia konkretnej metryki pasującej do sposobu niezbalansowania.

Dane podzieliłem na zbiór treningowy (fit) i walidacyjny (val) w proporcji 80:20, z zachowaniem proporcji klas. Do optymalizacji hiperparametrów zastosowałem StratifiedKFold z 10 foldami.

\section{Optymalizacja hiperparametrów drzewa decyzyjnego}

Przeprowadziłem optymalizację hiperparametrów klasyfikatora Decision Tree z wykorzystaniem przeszukiwania siatki parametrów i walidacji krzyżowej.

Analizowałem wpływ następujących parametrów:
\begin{itemize}
    \item criterion - sposób oceny jakości podziału (gini/entropy),
    \item max\_depth - maksymalna głębokość drzewa,
    \item max\_leaf\_nodes - maksymalna liczba liści.
\end{itemize}
Najlepszy uzyskany model wykorzystywał kryterium entropy, maksymalną głębokość równą 7 oraz maksymalną liczbę liści równą 18. Uzyskany wynik ROC AUC podczas walidacji krzyżowej wyniósł 0.7412.

\section{Podsumowanie i wnioski}

Wytrenowany model uzyskał ROC AUC równy 0,7759 na zbiorze treningowym, 0,7538 na zbiorze walidacyjnym, i 0,7716 na całości.

Przeprowadzone eksperymenty pozwoliły określić wpływ złożoności drzewa na jakość predykcji oraz wybrać zestaw hiperparametrów stanowiący punkt odniesienia dla kolejnego etapu projektu.

% ------------------------------- %
% ---  KONIEC TREŚCI RAPORTU  --- %
% ------------------------------- %

\end{document}
